% !TEX program = xelatex
\documentclass[UTF8,fontset=fandol,a4paper,11pt]{ctexart}
\usepackage[margin=21mm,top=23mm,bottom=22mm,headheight=16pt]{geometry}
\usepackage{amsmath,amssymb,mathtools,bm}
\usepackage{tikz}
\usetikzlibrary{arrows.meta,calc,angles,quotes,patterns}
\usepackage{xcolor,booktabs,array,longtable,enumitem,fancyhdr,needspace}
\usepackage[unicode,hidelinks]{hyperref}
\definecolor{ink}{HTML}{24343C}
\definecolor{accent}{HTML}{166A70}
\definecolor{muted}{HTML}{65747A}
\definecolor{rulegray}{HTML}{D4DDE0}
\definecolor{warm}{HTML}{966037}
\color{ink}
\setlength{\parindent}{2em}
\setlength{\parskip}{4pt}
\linespread{1.12}
\AtBeginDocument{\setlength{\abovedisplayskip}{6pt plus 2pt minus 1pt}
\setlength{\belowdisplayskip}{6pt plus 2pt minus 1pt}
\setlength{\abovedisplayshortskip}{0pt plus 1pt}
\setlength{\belowdisplayshortskip}{4pt plus 1pt}}
\setlist[enumerate]{label=(\arabic*),leftmargin=2.3em,itemsep=3pt,topsep=4pt}
\setlist[itemize]{leftmargin=2em,itemsep=2pt,topsep=4pt}
\allowdisplaybreaks[1]
\setlength{\emergencystretch}{2em}
\setcounter{secnumdepth}{0}
\providecommand{\tightlist}{\setlength{\itemsep}{0pt}\setlength{\parskip}{0pt}}
\ctexset{section={format=\Large\bfseries\color{accent},beforeskip=10pt,afterskip=9pt},
subsection={format=\normalsize\bfseries\color{accent},beforeskip=10pt,afterskip=4pt},
subsubsection={format=\normalsize\bfseries,beforeskip=8pt,afterskip=3pt}}
\pagestyle{fancy}
\fancyhf{}
\fancyhead[L]{\small\color{muted}物理学难题集萃}
% 页眉仅保留书名；版本名称只在封面显示。
\fancyfoot[L]{\small\color{muted}\chapterlabel}
\fancyfoot[R]{\small\thepage}
\renewcommand{\headrulewidth}{0.4pt}
\renewcommand{\footrulewidth}{0pt}
\newcommand{\dd}{\mathop{}\!\mathrm{d}}
\newcommand{\vect}[1]{\bm{#1}}
\newcommand{\tagline}[1]{\noindent{\small\color{muted}#1}\par}
\newcommand{\result}[1]{\begin{equation*}\boxed{#1}\end{equation*}}
\newcommand{\figcaption}[1]{\par\vspace{2pt}{\small\color{muted}#1}\par}
\newcommand{\teach}[2]{\Needspace{4\baselineskip}\subsection*{讲解要点与追问}#1\par\smallskip\noindent\textbf{追问：}#2\par}
\newcommand{\workarea}{\par\vspace{7mm}\noindent{\small\color{muted}作答区}\par\vspace{4mm}
\noindent\textcolor{rulegray}{\rule{\linewidth}{0.3pt}}\par\vspace{16mm}
\noindent\textcolor{rulegray}{\rule{\linewidth}{0.3pt}}\par\vspace{16mm}
\noindent\textcolor{rulegray}{\rule{\linewidth}{0.3pt}}\par\vspace{16mm}
\noindent\textcolor{rulegray}{\rule{\linewidth}{0.3pt}}}
\tikzset{axis/.style={-{Stealth[length=2mm]},thin,color=muted},
motion/.style={-{Stealth[length=2.3mm]},thick,color=accent},
force/.style={-{Stealth[length=2.3mm]},thick,color=warm},
traj/.style={thick,color=accent},guide/.style={dashed,thin,color=muted},
object/.style={thick,color=ink},every node/.style={font=\small}}\widowpenalty=10000\clubpenalty=10000\displaywidowpenalty=10000\newcommand{\chapterlabel}{力学试题（二）}







\hypersetup{pdftitle={力学试题二 力学试题（二）·学生版},pdfauthor={}}

\begin{document}\pdfbookmark[0]{力学试题二 力学试题（二）}{cover}

\thispagestyle{empty}\vspace*{22mm}

\noindent{\Huge\bfseries 物理学难题集萃}\par\vspace{10mm}

\noindent{\LARGE 力学试题二}\par\vspace{6mm}

\noindent{\LARGE 力学试题（二）}\par\vspace{12mm}

\noindent{\Large\color{accent}学生版}\par\vspace{9mm}

\noindent\textcolor{accent}{\rule{\linewidth}{1pt}}\par\vspace{9mm}

\noindent 第1--6题\par\vfill\clearpage

\pdfbookmark[1]{题目目录}{contents}\section*{题目目录}

\begin{longtable}{@{}r p{113mm} r@{}}\toprule 题号 & 题目 & 页码\\\midrule

\hyperref[q:1]{1} & \hyperref[q:1]{移动圆环与带环套的地面轨迹} & \pageref{q:1} \\

\hyperref[q:2]{2} & \hyperref[q:2]{三次喷气中的能量、速度矢量与落地点} & \pageref{q:2} \\

\hyperref[q:3]{3} & \hyperref[q:3]{绳越过凸圆弧时的张力与脱离} & \pageref{q:3} \\

\hyperref[q:4]{4} & \hyperref[q:4]{两细杆端部弹性碰撞与接触点速度} & \pageref{q:4} \\

\hyperref[q:5]{5} & \hyperref[q:5]{对称铰接双杆剪绳后的下落} & \pageref{q:5} \\

\hyperref[q:6]{6} & \hyperref[q:6]{双线悬圆柱的平移、复摆与偏航} & \pageref{q:6} \\

\bottomrule\end{longtable}

\clearpage

\pdfbookmark[2]{第 1 题 移动圆环与带环套的地面轨迹}{q1}

\section*{第 1 题\quad 移动圆环与带环套的地面轨迹}\label{q:1}


光滑水平面上放着均匀圆环，质量为$M$、半径为$R$。质量为$m<M$的质点可沿环的光滑内壁滑动。开始时圆环静止，环心在$O=(0,0)$，质点位于$(0,R)$，初速度为$v_0\hat x$。运动始终在水平面内进行。
\begin{enumerate}
\item 证明质点不会离开内壁。
\item 求质点相对地面的运动方程，并画出轨迹。
\item 求轨迹在最高处、最低处和竖直切线处的曲率半径；说明一周期内各类点出现的位置。
\end{enumerate}


\begin{center}\begin{minipage}{.98\linewidth}\centering
\begin{tikzpicture}[scale=.9,transform shape]
\draw[axis] (0.00000,0.00000)--(4.25000,0.00000);\draw[axis] (0.00000,-3.50000)--(0.00000,4.25000);\node[below right] at (4.25000,0.00000) {$x$};\node[above right] at (0.00000,4.25000) {$y$};\draw[thin,color=muted] (-3.50000,0.00000)--(0.00000,0.00000);\draw[object] (0.00000,0.00000) circle(3.00000);\draw[object] (0.00000,0.00000) circle(3.10000);\fill (0.00000,3.00000) circle(.035);\node[above left] at (0.00000,3.00000) {$m$};\fill (0.00000,0.00000) circle(.035);\node[below left] at (0.00000,0.00000) {$O$};\draw[motion] (0.00000,3.00000)--(2.00000,3.00000);\node[above] at (1.00000,3.10000) {$v_0$};\draw[<->,thin,color=muted] (0.00000,-0.60000)--(3.00000,-0.60000);\node[below] at (1.50000,-0.60000) {$R$};\node[below] at (-2.00000,-3.60000) {圆环 $M$，水平面};
\end{tikzpicture}
\figcaption{图力试2.1\quad 初始质点在环内壁顶端，初速沿+x，运动在水平面内。}
\end{minipage}\end{center}

\workarea

\clearpage

\pdfbookmark[2]{第 2 题 三次喷气中的能量、速度矢量与落地点}{q2}

\section*{第 2 题\quad 三次喷气中的能量、速度矢量与落地点}\label{q:2}


飞船起初在半径$r_0$的圆轨道上绕地球运动。每次喷气持续时间极短，喷气过程中位置不变，各阶段把飞船质量视为常量。将单位质量的冲量$\Delta\boldsymbol v=\Delta\boldsymbol p/m$称为比冲量，其大小单位为速度。地球质量为$M_E$、半径为$R_E$，$r_1>r_0>R_E$，忽略空气阻力。
\begin{enumerate}
\item 第一次喷气使飞船刚能逃逸，求最小比冲量及排气方向。
\item 第一次取上述最小值。飞船沿上升轨迹到达$r_1$处，第二次喷气使它进入半径$r_1$的同向圆轨道。求比冲量及排气方向。
\item 飞船在该圆轨道的$A$点第三次喷气，希望在对径方向的地表点$B$与地表相切。分别讨论向正前方排气和沿径向向外排气；检查两者是否都能到达指定的$B$。若只要求在某地表点相切而不限定方位，求两种比冲量，并在$r_0=2R_E$时比较。
\end{enumerate}
各次排气方向指相对飞船的喷气方向，与飞船获得的冲量方向相反。“与地表相切”指飞船到地表时的瞬时速度沿当地切线，落地速率由轨道确定。


\begin{center}\begin{minipage}{.98\linewidth}\centering
\begin{tikzpicture}[scale=.9,transform shape]
\draw[object] (0.00000,0.00000) circle(1.00000);\draw[guide] (0.00000,0.00000) circle(2.00000);\draw[guide] (0.00000,0.00000) circle(4.00000);\fill (0.00000,0.00000) circle(.035);\node[below] at (0.00000,0.00000) {$O$};\fill (4.00000,0.00000) circle(.035);\node[right] at (4.00000,0.00000) {$A$};\fill (-1.00000,0.00000) circle(.035);\node[left] at (-1.00000,0.00000) {$B$};\draw[guide] (-1.00000,0.00000)--(4.00000,0.00000);\node[above] at (0.00000,1.00000) {$R_E$};\node[above left] at (-2.40000,0.40000) {$r_0$};\node[above] at (0.00000,4.00000) {$r_1$};\draw[motion] (-2.00000,0.00000)--(-2.00000,-1.30000);\node[left] at (-2.10000,-1.00000) {$v_0$};
\end{tikzpicture}
\figcaption{图力试2.2\quad 三条同心尺度及第三次喷气的指定对径地表点B。}
\end{minipage}\end{center}

\par\vspace{5mm}\noindent{\small\color{muted}作答区}\par\vspace{3mm}\noindent\textcolor{rulegray}{\rule{\linewidth}{0.3pt}}\par\vspace{12mm}\noindent\textcolor{rulegray}{\rule{\linewidth}{0.3pt}}

\clearpage

\pdfbookmark[2]{第 3 题 绳越过凸圆弧时的张力与脱离}{q3}

\section*{第 3 题\quad 绳越过凸圆弧时的张力与脱离}\label{q:3}


一条均匀柔绳总长为$l$、线密度为$\lambda$，放在光滑水平面上，右端位于与下降斜面相接的凸边$B$。斜面倾角为$\theta$，$0<\theta<\pi/2$；连接处用半径$r\ll l$的小圆弧平滑过渡。绳可自由弯曲且不可伸长。

让绳端因微小扰动先下滑一小段，随后在重力下运动。以$x$表示尚在水平面上的绳长。求下滑过程中的绳张力、凸边对绳的支持力，以及绳开始脱离凸边时的位置。先求$r/l\to0$的领先阶结果，并说明有限$r$时脱离条件如何修正。圆弧没有黏附能力，只能向外推绳。


\begin{center}\begin{minipage}{.98\linewidth}\centering
\begin{tikzpicture}[scale=.9,transform shape]
\draw[object] (-4.00000,0.00000)--(0.00000,0.00000);\draw[object] plot coordinates {(0.00000,-0.00000) (0.00740,-0.00003) (0.01481,-0.00012) (0.02221,-0.00027) (0.02961,-0.00049) (0.03701,-0.00076) (0.04441,-0.00110) (0.05181,-0.00149) (0.05920,-0.00195) (0.06658,-0.00247) (0.07397,-0.00304) (0.08134,-0.00368) (0.08872,-0.00438) (0.09608,-0.00514) (0.10344,-0.00596) (0.11079,-0.00685) (0.11814,-0.00779) (0.12548,-0.00879) (0.13280,-0.00985) (0.14012,-0.01098) (0.14743,-0.01216) (0.15473,-0.01340) (0.16202,-0.01470) (0.16930,-0.01607) (0.17657,-0.01749) (0.18382,-0.01897) (0.19107,-0.02052) (0.19830,-0.02212) (0.20551,-0.02378) (0.21271,-0.02550) (0.21990,-0.02728) (0.22708,-0.02912) (0.23423,-0.03102) (0.24138,-0.03297) (0.24850,-0.03499) (0.25561,-0.03706) (0.26270,-0.03919) (0.26977,-0.04138) (0.27683,-0.04363) (0.28387,-0.04594) (0.29088,-0.04830) (0.29788,-0.05073) (0.30486,-0.05321) (0.31182,-0.05574) (0.31875,-0.05834) (0.32567,-0.06099) (0.33256,-0.06370) (0.33943,-0.06646) (0.34627,-0.06928) (0.35310,-0.07216) (0.35990,-0.07509) (0.36667,-0.07808) (0.37342,-0.08113) (0.38015,-0.08423) (0.38685,-0.08738) (0.39352,-0.09059) (0.40017,-0.09386) (0.40678,-0.09718) (0.41338,-0.10055) (0.41994,-0.10398) (0.42648,-0.10746) (0.43298,-0.11100) (0.43946,-0.11459) (0.44591,-0.11823) (0.45232,-0.12192) (0.45871,-0.12567) (0.46507,-0.12947) (0.47139,-0.13332) (0.47768,-0.13723) (0.48394,-0.14118) (0.49017,-0.14519) (0.49636,-0.14925) (0.50252,-0.15336) (0.50865,-0.15752) (0.51474,-0.16173) (0.52080,-0.16599) (0.52682,-0.17030) (0.53280,-0.17466) (0.53875,-0.17907) (0.54467,-0.18352)};\draw[object] (0.54467,-0.18352)--(3.72900,-2.60427);\draw[very thick,color=accent] (-3.00000,0.00000)--(0.00000,0.00000);\draw[traj] plot coordinates {(0.00000,-0.00000) (0.00740,-0.00003) (0.01481,-0.00012) (0.02221,-0.00027) (0.02961,-0.00049) (0.03701,-0.00076) (0.04441,-0.00110) (0.05181,-0.00149) (0.05920,-0.00195) (0.06658,-0.00247) (0.07397,-0.00304) (0.08134,-0.00368) (0.08872,-0.00438) (0.09608,-0.00514) (0.10344,-0.00596) (0.11079,-0.00685) (0.11814,-0.00779) (0.12548,-0.00879) (0.13280,-0.00985) (0.14012,-0.01098) (0.14743,-0.01216) (0.15473,-0.01340) (0.16202,-0.01470) (0.16930,-0.01607) (0.17657,-0.01749) (0.18382,-0.01897) (0.19107,-0.02052) (0.19830,-0.02212) (0.20551,-0.02378) (0.21271,-0.02550) (0.21990,-0.02728) (0.22708,-0.02912) (0.23423,-0.03102) (0.24138,-0.03297) (0.24850,-0.03499) (0.25561,-0.03706) (0.26270,-0.03919) (0.26977,-0.04138) (0.27683,-0.04363) (0.28387,-0.04594) (0.29088,-0.04830) (0.29788,-0.05073) (0.30486,-0.05321) (0.31182,-0.05574) (0.31875,-0.05834) (0.32567,-0.06099) (0.33256,-0.06370) (0.33943,-0.06646) (0.34627,-0.06928) (0.35310,-0.07216) (0.35990,-0.07509) (0.36667,-0.07808) (0.37342,-0.08113) (0.38015,-0.08423) (0.38685,-0.08738) (0.39352,-0.09059) (0.40017,-0.09386) (0.40678,-0.09718) (0.41338,-0.10055) (0.41994,-0.10398) (0.42648,-0.10746) (0.43298,-0.11100) (0.43946,-0.11459) (0.44591,-0.11823) (0.45232,-0.12192) (0.45871,-0.12567) (0.46507,-0.12947) (0.47139,-0.13332) (0.47768,-0.13723) (0.48394,-0.14118) (0.49017,-0.14519) (0.49636,-0.14925) (0.50252,-0.15336) (0.50865,-0.15752) (0.51474,-0.16173) (0.52080,-0.16599) (0.52682,-0.17030) (0.53280,-0.17466) (0.53875,-0.17907) (0.54467,-0.18352)};\draw[traj] (0.54467,-0.18352)--(2.77370,-1.87805);\fill (0.00000,0.00000) circle(.035);\node[above left] at (0.00000,0.00000) {$B$};\fill (0.00000,-0.90000) circle(.035);\node[below left] at (0.00000,-0.90000) {曲率中心};\draw[guide] (0.00000,-0.90000)--(0.54467,-0.18352);\node[right] at (0.20000,-0.55000) {$r$};\draw[<->,thin,color=muted] (-3.00000,1.00000)--(0.00000,1.00000);\node[below] at (-1.50000,1.00000) {$x$};\node[above] at (2.70000,-0.30000) {下降斜面 $\theta$};\node[above] at (0.00000,2.10000) {圆角放大；$r\ll l$};
\end{tikzpicture}
\figcaption{图力试2.3\quad 凸边半径很小，几何图将圆角放大显示。}
\end{minipage}\end{center}

\workarea

\clearpage

\pdfbookmark[2]{第 4 题 两细杆端部弹性碰撞与接触点速度}{q4}

\section*{第 4 题\quad 两细杆端部弹性碰撞与接触点速度}\label{q:4}


光滑水平桌面上，两根均匀细杆$A,B$平行于$x$轴，长度均为$l$，质量分别为$m_A,m_B$。杆$A$静止，位于$x$轴的$(-l+\epsilon,\epsilon)$区域，$0<\epsilon\ll l$。杆$B$位于其下方，横坐标覆盖$(0,l)$，以$v_0\hat y$作无转动平移。

当杆$B$向上到达$x$轴时，其左端与杆$A$右端发生光滑、单点、完全弹性碰撞。取$\epsilon/l\to0$的极限，接触法向沿$y$；忽略杆厚和接触时间内的位移。
\begin{enumerate}
\item 求碰后两杆质心的带符号速度$v_A,v_B$及角速度$\omega_A,\omega_B$，其中向$+y$、逆时针分别为正。
\item 验证
\[
v_0=\left(v_A+\frac l2\omega_A\right)
-\left(v_B-\frac l2\omega_B\right),
\]
并说明它在本题中的物理依据和适用条件。
\end{enumerate}


\begin{center}\begin{minipage}{.98\linewidth}\centering
\begin{tikzpicture}[scale=.9,transform shape]
\draw[axis] (0.00000,0.00000)--(5.25000,0.00000);\draw[axis] (0.00000,-2.20000)--(0.00000,2.25000);\node[below right] at (5.25000,0.00000) {$x$};\node[above right] at (0.00000,2.25000) {$y$};\draw[very thick,color=accent] (-4.00000,0.00000)--(0.00000,0.00000);\draw[very thick,color=accent] (0.00000,-1.50000)--(4.00000,-1.50000);\node[below] at (-2.00000,-0.10000) {杆 $A,m_A$};\node[below] at (2.00000,-1.60000) {杆 $B,m_B$};\draw[motion] (2.00000,-1.40000)--(2.00000,-0.30000);\node[right] at (2.10000,-0.70000) {$v_0$};\fill (0.00000,0.00000) circle(.035);\node[below left] at (0.00000,0.00000) {$O$};
\end{tikzpicture}
\figcaption{图力试2.4\quad B杆向+y平移，端部将与静止A杆接触，取小重叠极限。}
\end{minipage}\end{center}

\workarea

\clearpage

\pdfbookmark[2]{第 5 题 对称铰接双杆剪绳后的下落}{q5}

\section*{第 5 题\quad 对称铰接双杆剪绳后的下落}\label{q:5}


两根相同均匀细杆各长$l$、质量各为$m$，上端以光滑铰链$C$连接，下端放在光滑水平地面上，成对称的人字形。起初每杆与水平面夹角$\theta_0=30^\circ$，用一根轻绳保持静止。突然剪断轻绳，双杆在竖直平面内运动，忽略阻力。
\begin{enumerate}
\item 求铰链$C$即将碰到地面时的速度。
\item 用积分式表示从剪绳到铰链碰地所需的时间，并检查两杆下端在这段时间内是否保持接地。
\end{enumerate}
只求第一次碰地之前的运动，不讨论碰地冲击。


\begin{center}\begin{minipage}{.98\linewidth}\centering
\begin{tikzpicture}[scale=.9,transform shape]
\draw[object] (-5.00000,0.00000)--(5.00000,0.00000);\draw[very thick,color=accent] (-3.46410,0.00000)--(0.00000,2.00000)--(3.46410,0.00000);\fill (0.00000,2.00000) circle(.035);\node[above] at (0.00000,2.00000) {$C$};\draw[guide] (-1.73205,1.00000)--(1.73205,1.00000);\node[above] at (0.00000,1.10000) {剪断轻绳};\node[above left] at (-2.60000,1.40000) {$l,m$};\node[above right] at (2.60000,1.40000) {$l,m$};\node[below] at (-2.90000,-0.10000) {$\theta_0$};\node[below] at (2.90000,-0.10000) {$\theta_0$};
\end{tikzpicture}
\figcaption{图力试2.5\quad 两杆对称，绳剪断后脚点向两侧张开，铰链竖直下降。}
\end{minipage}\end{center}

\workarea

\clearpage

\pdfbookmark[2]{第 6 题 双线悬圆柱的平移、复摆与偏航}{q6}

\section*{第 6 题\quad 双线悬圆柱的平移、复摆与偏航}\label{q:6}


均匀实心圆柱质量为$m$、半径为$R$、轴长为$L$，由两根等长为$l$的轻细线对称悬挂。两线沿柱轴方向相距$D$，$0<D\le L$，下端固定在圆柱上表面；上端固定点为$O,O'$。圆柱质心$C$位于两线中间正下方。

求下列小振幅运动的周期，并明确各自的约束：
\begin{enumerate}
\item 两线在包含柱轴的同一铅垂平面内同角摆动，圆柱保持平行平移。
\item 把圆柱与悬线作为绕$OO'$轴摆动的一自由度受限复摆，求周期；并判断仅用柔线自由悬挂时，这种整体绕轴运动是否自动成立。
\item 圆柱在水平面内绕通过质心的铅垂轴小角偏航，求周期，并由薄圆片积分求所需转动惯量。
\end{enumerate}
忽略线的质量、伸长和阻力。受限复摆模式的导向不作功，其余运动由张紧柔线实现。


\begin{center}\begin{minipage}{.98\linewidth}\centering
\begin{tikzpicture}[scale=.9,transform shape]
\draw[object] plot coordinates {(-2.78000,-3.00000) (-2.78044,-2.96195) (-2.78177,-2.92404) (-2.78398,-2.88645) (-2.78705,-2.84931) (-2.79098,-2.81278) (-2.79576,-2.77700) (-2.80136,-2.74212) (-2.80775,-2.70828) (-2.81492,-2.67562) (-2.82284,-2.64426) (-2.83147,-2.61433) (-2.84078,-2.58595) (-2.85073,-2.55924) (-2.86128,-2.53431) (-2.87239,-2.51125) (-2.88401,-2.49016) (-2.89610,-2.47113) (-2.90861,-2.45422) (-2.92149,-2.43951) (-2.93468,-2.42706) (-2.94813,-2.41691) (-2.96180,-2.40912) (-2.97562,-2.40370) (-2.98953,-2.40068) (-3.00349,-2.40008) (-3.01743,-2.40189) (-3.03131,-2.40611) (-3.04506,-2.41272) (-3.05862,-2.42169) (-3.07195,-2.43300) (-3.08500,-2.44659) (-3.09769,-2.46240) (-3.11000,-2.48038) (-3.12186,-2.50046) (-3.13323,-2.52254) (-3.14407,-2.54655) (-3.15432,-2.57238) (-3.16396,-2.59994) (-3.17293,-2.62910) (-3.18121,-2.65976) (-3.18876,-2.69179) (-3.19554,-2.72506) (-3.20154,-2.75944) (-3.20673,-2.79479) (-3.21109,-2.83096) (-3.21459,-2.86781) (-3.21724,-2.90520) (-3.21900,-2.94297) (-3.21989,-2.98096) (-3.21989,-3.01904) (-3.21900,-3.05703) (-3.21724,-3.09480) (-3.21459,-3.13219) (-3.21109,-3.16904) (-3.20673,-3.20521) (-3.20154,-3.24056) (-3.19554,-3.27494) (-3.18876,-3.30821) (-3.18121,-3.34024) (-3.17293,-3.37090) (-3.16396,-3.40006) (-3.15432,-3.42762) (-3.14407,-3.45345) (-3.13323,-3.47746) (-3.12186,-3.49954) (-3.11000,-3.51962) (-3.09769,-3.53760) (-3.08500,-3.55341) (-3.07195,-3.56700) (-3.05862,-3.57831) (-3.04506,-3.58728) (-3.03131,-3.59389) (-3.01743,-3.59811) (-3.00349,-3.59992) (-2.98953,-3.59932) (-2.97562,-3.59630) (-2.96180,-3.59088) (-2.94813,-3.58309) (-2.93468,-3.57294) (-2.92149,-3.56049) (-2.90861,-3.54578) (-2.89610,-3.52887) (-2.88401,-3.50984) (-2.87239,-3.48875) (-2.86128,-3.46569) (-2.85073,-3.44076) (-2.84078,-3.41405) (-2.83147,-3.38567) (-2.82284,-3.35574) (-2.81492,-3.32438) (-2.80775,-3.29172) (-2.80136,-3.25788) (-2.79576,-3.22300) (-2.79098,-3.18722) (-2.78705,-3.15069) (-2.78398,-3.11355) (-2.78177,-3.07596) (-2.78044,-3.03805) (-2.78000,-3.00000)};\draw[object] plot coordinates {(3.22000,-3.00000) (3.21956,-2.96195) (3.21823,-2.92404) (3.21602,-2.88645) (3.21295,-2.84931) (3.20902,-2.81278) (3.20424,-2.77700) (3.19864,-2.74212) (3.19225,-2.70828) (3.18508,-2.67562) (3.17716,-2.64426) (3.16853,-2.61433) (3.15922,-2.58595) (3.14927,-2.55924) (3.13872,-2.53431) (3.12761,-2.51125) (3.11599,-2.49016) (3.10390,-2.47113) (3.09139,-2.45422) (3.07851,-2.43951) (3.06532,-2.42706) (3.05187,-2.41691) (3.03820,-2.40912) (3.02438,-2.40370) (3.01047,-2.40068) (2.99651,-2.40008) (2.98257,-2.40189) (2.96869,-2.40611) (2.95494,-2.41272) (2.94138,-2.42169) (2.92805,-2.43300) (2.91500,-2.44659) (2.90231,-2.46240) (2.89000,-2.48038) (2.87814,-2.50046) (2.86677,-2.52254) (2.85593,-2.54655) (2.84568,-2.57238) (2.83604,-2.59994) (2.82707,-2.62910) (2.81879,-2.65976) (2.81124,-2.69179) (2.80446,-2.72506) (2.79846,-2.75944) (2.79327,-2.79479) (2.78891,-2.83096) (2.78541,-2.86781) (2.78276,-2.90520) (2.78100,-2.94297) (2.78011,-2.98096) (2.78011,-3.01904) (2.78100,-3.05703) (2.78276,-3.09480) (2.78541,-3.13219) (2.78891,-3.16904) (2.79327,-3.20521) (2.79846,-3.24056) (2.80446,-3.27494) (2.81124,-3.30821) (2.81879,-3.34024) (2.82707,-3.37090) (2.83604,-3.40006) (2.84568,-3.42762) (2.85593,-3.45345) (2.86677,-3.47746) (2.87814,-3.49954) (2.89000,-3.51962) (2.90231,-3.53760) (2.91500,-3.55341) (2.92805,-3.56700) (2.94138,-3.57831) (2.95494,-3.58728) (2.96869,-3.59389) (2.98257,-3.59811) (2.99651,-3.59992) (3.01047,-3.59932) (3.02438,-3.59630) (3.03820,-3.59088) (3.05187,-3.58309) (3.06532,-3.57294) (3.07851,-3.56049) (3.09139,-3.54578) (3.10390,-3.52887) (3.11599,-3.50984) (3.12761,-3.48875) (3.13872,-3.46569) (3.14927,-3.44076) (3.15922,-3.41405) (3.16853,-3.38567) (3.17716,-3.35574) (3.18508,-3.32438) (3.19225,-3.29172) (3.19864,-3.25788) (3.20424,-3.22300) (3.20902,-3.18722) (3.21295,-3.15069) (3.21602,-3.11355) (3.21823,-3.07596) (3.21956,-3.03805) (3.22000,-3.00000)};\draw[object] (-3.00000,-2.40000)--(3.00000,-2.40000);\draw[object] (-3.00000,-3.60000)--(3.00000,-3.60000);\draw[object] (-2.30000,0.00000)--(2.30000,0.00000);\fill (-1.70000,0.00000) circle(.035);\node[above] at (-1.70000,0.00000) {$O$};\draw[object] (-1.70000,0.00000)--(-1.70000,-2.40000);\fill (1.70000,0.00000) circle(.035);\node[above] at (1.70000,0.00000) {$O'$};\draw[object] (1.70000,0.00000)--(1.70000,-2.40000);\fill (0.00000,-3.00000) circle(.035);\node[below] at (0.00000,-3.00000) {$C$};\draw[<->,thin,color=muted] (-1.70000,0.90000)--(1.70000,0.90000);\node[below] at (0.00000,0.90000) {$D$};\draw[<->,thin,color=muted] (4.00000,0.00000)--(4.00000,-2.40000);\node[right] at (4.00000,-1.20000) {$l$};\draw[<->,thin,color=muted] (-4.00000,-3.00000)--(-4.00000,-2.40000);\node[left] at (-4.00000,-2.70000) {$R$};\draw[<->,thin,color=muted] (-3.00000,-4.50000)--(3.00000,-4.50000);\node[below] at (0.00000,-4.50000) {$L$};
\end{tikzpicture}
\figcaption{图力试2.6\quad 柱轴方向两吊点相距D，线长l，质心低于吊点R。}
\end{minipage}\end{center}

\workarea

\end{document}
